% !TEX program = pdflatex
\documentclass[11pt,a4paper]{article}

% =========================================================================
% BOSQUEJO DEL CIERRE DEL PROYECTO
% Proyecto: CasaScore (AlquilaPuno)
% Gestión de Proyectos TICS (EST315) -- Semestre 2026-II -- Team Sobrevivientes
% =========================================================================

% --- PAQUETES DE CONFIGURACIÓN ---
\usepackage[utf8]{inputenc}
\usepackage[T1]{fontenc}
\usepackage[spanish,es-tabla]{babel}
\usepackage{geometry}
\geometry{
    top=3.6cm,
    bottom=2.3cm,
    left=2.2cm,
    right=2.2cm,
    % PARTE MODIFICABLE: headheight debe ser mayor que la altura del logo más alto (2.56cm = 73pt).
    headheight=80pt,
    headsep=15pt,
    footskip=20pt
}

\usepackage{graphicx}
\usepackage{fancyhdr}
\usepackage[table]{xcolor}
\usepackage{titlesec}
\usepackage{booktabs}
\usepackage{tabularx}
\usepackage{array}
\usepackage{amsmath}
\usepackage{enumitem}
\usepackage{microtype}
\usepackage{amssymb}
% --- RECURSOS VISUALES: GRÁFICOS Y CÓDIGO DE COLORES ---
\usepackage{tikz}
\usetikzlibrary{positioning,calc}
\definecolor{GrisFila}{RGB}{246, 248, 251}
\definecolor{CuadCerca}{RGB}{150, 30, 30}
\definecolor{CuadSatisf}{RGB}{180, 130, 20}
\definecolor{CuadInfo}{RGB}{10, 45, 90}
\definecolor{CuadMonit}{RGB}{110, 115, 125}
\newcommand{\etiqueta}[2]{{\setlength{\fboxsep}{2pt}\colorbox{#1}{\textcolor{white}{\scriptsize\bfseries #2}}}}
\newcommand{\cerca}{\etiqueta{CuadCerca}{Gestionar de cerca}}
\newcommand{\satisf}{\etiqueta{CuadSatisf}{Mantener satisfecho}}
\newcommand{\informado}{\etiqueta{CuadInfo}{Mantener informado}}
\newcommand{\monit}{\etiqueta{CuadMonit}{Monitorear}}
\usepackage{needspace}
\usepackage{longtable}
\definecolor{Rama1}{RGB}{10, 45, 90}
\definecolor{Rama2}{RGB}{40, 95, 150}
\definecolor{Rama3}{RGB}{180, 130, 20}
\definecolor{Rama4}{RGB}{30, 110, 60}
\definecolor{Rama5}{RGB}{110, 60, 140}
\definecolor{RaciR}{RGB}{150, 30, 30}
\definecolor{RaciA}{RGB}{10, 45, 90}
\definecolor{RaciC}{RGB}{180, 130, 20}
\definecolor{RaciI}{RGB}{140, 145, 155}
\definecolor{NivAlto}{RGB}{30, 110, 60}
\definecolor{NivMedio}{RGB}{180, 130, 20}
\definecolor{NivBajo}{RGB}{150, 30, 30}
\usetikzlibrary{patterns}
\definecolor{SevBaja}{RGB}{30, 110, 60}
\definecolor{SevMedia}{RGB}{180, 130, 20}
\definecolor{SevAlta}{RGB}{150, 30, 30}
\usetikzlibrary{arrows.meta,shapes.geometric,fit,backgrounds}
\usepackage{listings}
\lstset{basicstyle=\ttfamily\scriptsize, frame=single, rulecolor=\color{AzulUNA!40}, backgroundcolor=\color{GrisFila}, columns=fullflexible, keepspaces=true, breaklines=true, xleftmargin=4pt, xrightmargin=4pt, literate={á}{{\'a}}1 {é}{{\'e}}1 {í}{{\'i}}1 {ó}{{\'o}}1 {ú}{{\'u}}1 {ñ}{{\~n}}1 {Á}{{\'A}}1}

\usepackage{pdflscape}
\usepackage{adjustbox}
\usepackage{afterpage}
\usetikzlibrary{shapes.multipart,shapes.misc,decorations.pathreplacing}
\usepackage{hyperref}
\AtBeginDocument{\shorthandoff{<>}}
% --- HOJAS HORIZONTALES: encabezado y pie girados para leerse con la hoja apaisada ---
\fancypagestyle{horizontal}{%
  \fancyhf{}\renewcommand{\headrulewidth}{0pt}\renewcommand{\footrulewidth}{0pt}%
  \fancyhead[C]{\begin{tikzpicture}[remember picture,overlay]
    \node[rotate=90, anchor=center] at ($(current page.west)+(1.25cm,0.75cm)$) {%
      \begin{tabular}{@{}c@{}}\begin{minipage}{24.6cm}
      \raisebox{-0.5\height}{\includegraphics[height=1.15cm,keepaspectratio]{logo_fiei.png}}\hfill
      \begin{tabular}{c}{\scriptsize\bfseries\textcolor{AzulUNA}{UNIVERSIDAD NACIONAL DEL ALTIPLANO DE PUNO}}\\[-1pt]
      {\scriptsize\bfseries\textcolor{AzulUNA}{Gestión de Proyectos TICS (EST315) \textbar\ Team Sobrevivientes}}\\[-1pt]
      {\footnotesize\bfseries\textcolor{DoradoFIEI}{BOSQUEJO DEL CIERRE DEL PROYECTO}}\end{tabular}\hfill
      \raisebox{-0.5\height}{\includegraphics[height=2.0cm,keepaspectratio]{logo_equipo.png}}
      \end{minipage}\end{tabular}};
    \node[rotate=90, anchor=center] at ($(current page.east)+(-1.0cm,0.75cm)$) {%
      \begin{minipage}{24.6cm}{\scriptsize\textcolor{GrisOscuro}{CasaScore (AlquilaPuno)}}\hfill{\small\thepage}\hfill{\scriptsize\textcolor{GrisOscuro}{Puno, Perú -- Semestre 2026-II}}\end{minipage}};
  \end{tikzpicture}}}
\newenvironment{hojah}[1]{\clearpage\begin{landscape}\thispagestyle{horizontal}\section{#1}}{\end{landscape}}
\newenvironment{grande}[1][13.6cm]{\par\noindent\begin{adjustbox}{width=\linewidth, max totalheight=#1, center}}{\end{adjustbox}\par}
\newcommand{\tile}[6]{\node[draw=#3, fill=#3!8, line width=1pt, rounded corners=4pt, minimum width=3.95cm, minimum height=2.0cm] at (#1,#2) {};
\node[font=\scriptsize\bfseries, text=GrisOscuro] at (#1,#2+0.6) {#4};
\node[font=\Large\bfseries, text=#3] at (#1,#2) {#5};
\node[font=\tiny, text=GrisOscuro] at (#1,#2-0.6) {#6};}
\tikzset{
 fl/.style={>=Stealth, font=\normalsize, node distance=0.9cm and 1.0cm},
 term/.style={rounded rectangle, draw=AzulUNA, fill=AzulUNA, text=white, font=\normalsize\bfseries, minimum height=0.9cm, inner xsep=10pt},
 proc/.style={rectangle, rounded corners=3pt, draw=AzulUNA, fill=AzulClaro, line width=0.9pt, align=center, minimum height=1.0cm, text width=#1, inner sep=4pt},
 proc/.default=3.4cm,
 dec/.style={diamond, aspect=2.1, draw=DoradoFIEI, fill=DoradoFIEI!12, line width=0.9pt, align=center, inner sep=1pt, font=\small\bfseries, text width=#1},
 dec/.default=2.4cm,
 dato/.style={trapezium, trapezium left angle=72, trapezium right angle=108, draw=Rama4, fill=Rama4!10, line width=0.9pt, align=center, text width=#1, minimum height=1.0cm},
 dato/.default=3.2cm,
 fle/.style={->, line width=1.1pt, AzulUNA},
 etq/.style={font=\small\bfseries, text=GrisOscuro, fill=white, inner sep=1.5pt},
 actor/.style={draw=#1, fill=#1!10, line width=1.2pt, rounded corners=6pt, align=center, font=\large},
}
\newcommand{\marca}{{\Large$\square$}\,{\scriptsize Sí}\ \ {\Large$\square$}\,{\scriptsize No}}
\newcommand{\caja}[3]{\node[draw=#1, fill=#1!10, line width=1pt, rounded corners=4pt, align=center, inner sep=5pt] #2 {#3};}


% Logos: funciona a 2 o 3 niveles de profundidad dentro de docs/latex/
\graphicspath{{../../assets/img/}{../../../assets/img/}}

% --- PALETA INSTITUCIONAL FIEI / UNA PUNO ---
\definecolor{AzulUNA}{RGB}{10, 45, 90}
\definecolor{DoradoFIEI}{RGB}{180, 130, 20}
\definecolor{GrisOscuro}{RGB}{50, 50, 50}
\definecolor{AzulClaro}{RGB}{228, 236, 246}

% --- CONFIGURACIÓN DE LISTAS ---
\setlist{topsep=2pt, itemsep=2pt, parsep=1pt}

% --- COLUMNAS DE TABLA ---
\newcolumntype{L}{>{\raggedright\arraybackslash}X}
\newcolumntype{P}[1]{>{\raggedright\arraybackslash}p{#1}}
\newcolumntype{C}[1]{>{\centering\arraybackslash}p{#1}}
\newcolumntype{Q}[1]{>{\raggedleft\arraybackslash}p{#1}}
\newcolumntype{R}{>{\raggedleft\arraybackslash}X}
\renewcommand{\arraystretch}{1.25}

% --- FORMATO DE ENLACES ---
\hypersetup{
    colorlinks=true,
    linkcolor=AzulUNA,
    urlcolor=AzulUNA,
    citecolor=AzulUNA,
    pdftitle={Bosquejo del Cierre del Proyecto -- CasaScore (AlquilaPuno)},
    pdfauthor={Team Sobrevivientes}
}

% --- FORMATO DE TÍTULOS ---
\titleformat{\section}{\color{AzulUNA}\normalfont\large\bfseries}{\thesection.}{0.7em}{}[\color{AzulUNA}\titlerule]
\titleformat{\subsection}{\color{AzulUNA}\normalfont\normalsize\bfseries}{\thesubsection.}{0.5em}{}
\titlespacing*{\section}{0pt}{9pt}{4pt}
\titlespacing*{\subsection}{0pt}{6pt}{2pt}

% --- ENCABEZADO Y PIE DE PÁGINA (fancyhdr) ---
\pagestyle{fancy}
\fancyhf{}

% ===== TAMAÑO DE LOS LOGOS (PARTE MODIFICABLE) =====
% logo_fiei.png  : height=1.45cm (lado izquierdo)
% logo_equipo.png: height=2.56cm (lado derecho; 25 % más grande porque la imagen tiene más borde blanco)
% Para cambiar el tamaño, modifique solo el valor height=... de cada \includegraphics.
% Si el logo más alto supera 2.6cm, aumente también headheight y top en la configuración de página.
\fancyhead[L]{%
    \includegraphics[height=1.45cm,keepaspectratio]{logo_fiei.png}%
}
\fancyhead[C]{%
    \scriptsize \textbf{\textcolor{AzulUNA}{UNIVERSIDAD NACIONAL DEL ALTIPLANO DE PUNO}}\\[2pt]
    \tiny \textcolor{GrisOscuro}{FACULTAD DE INGENIERÍA ESTADÍSTICA E INFORMÁTICA}\\[2pt]
    \scriptsize \textbf{\textcolor{AzulUNA}{Gestión de Proyectos TICS (EST315) \textbar\ Team Sobrevivientes}}\\[2pt]
    \footnotesize \textbf{\textcolor{DoradoFIEI}{BOSQUEJO DEL CIERRE DEL PROYECTO}}%
}
\fancyhead[R]{%
    \includegraphics[height=2.56cm,keepaspectratio]{logo_equipo.png}%
}

\renewcommand{\headrulewidth}{0pt}
\renewcommand{\headrule}{}

\fancyfoot[L]{\scriptsize \textcolor{GrisOscuro}{CasaScore (AlquilaPuno)}}
\fancyfoot[C]{\small \thepage}
\fancyfoot[R]{\scriptsize \textcolor{GrisOscuro}{Puno, Perú -- Semestre 2026-II}}
\renewcommand{\footrulewidth}{0.4pt}

\begin{document}
\begin{center}
    {\Large \textbf{\textcolor{AzulUNA}{BOSQUEJO DEL CIERRE DEL PROYECTO}}}\\[3pt]
    {\normalsize \textbf{CasaScore (AlquilaPuno)}}\\[4pt]
    \footnotesize \textbf{Fecha:} viernes 13/11/2026 (semana 12) \quad \textbar\quad \textbf{Versión:} 0.2 (primer bosquejo completo) \quad \textbar\quad \textbf{Fase:} Cierre
\end{center}

\vspace{0.1cm}

\section{Datos Generales}
{\small
\noindent\begin{tabularx}{\textwidth}{@{}P{4.0cm} L@{}}
\toprule
\textbf{Proyecto:} & CasaScore (aplicación AlquilaPuno) \\
\textbf{Responsable:} & Clyde Neil Paricahua Pari -- Project Manager \\
\textbf{Participan:} & Ruth K. Apaza Solis, Eliseo Tarqui Ajahuana y Viviana M. Quisbert Quispe \\
\textbf{Proceso PMBOK:} & Cerrar el proyecto o fase \\
\textbf{Insumos:} & Plan de Dirección, informe de desempeño, registro de riesgos 1.1, registro de cambios y estado del tablero \\
\textbf{Se completa:} & Versión final con los resultados del piloto (27/11) y del último corte de valor ganado (11/12) \\
\bottomrule
\end{tabularx}
}

\section{Qué se Cierra y Qué Continúa}
CasaScore es un producto pensado para tres años o más (2027 a 2029). En diciembre de 2026 se cierra la \textbf{fase 1}, el proyecto del semestre, que entrega un MVP funcionando en la nube. El producto no se apaga: pasa a operación y a nuevas fases.

\vspace{0.1cm}
{\small
\rowcolors{2}{white}{GrisFila}
\noindent\begin{tabularx}{\textwidth}{@{}P{3.6cm} L L@{}}
\toprule
\rowcolor{AzulClaro}
& \textbf{Se cierra en la fase 1 (2026)} & \textbf{Continúa después del semestre} \\
\midrule
Alcance & MVP con las 9 funciones básicas (registro, publicación, punto de referencia, prioridades, puntaje, ranking, explicación, perfil y nube), piloto con clientes y expediente de gestión. & Mapa de cuartos, confirmación de alquiler, calificaciones, reportes, cobros dentro de la app y nuevas ciudades. \\
Producto & MVP publicado en la nube y probado con cuartos reales de compañeros. & Operación continua, mejoras por versión y app en Google Play. \\
Equipo & Se liberan las 8 horas semanales del semestre. & Equipo de operación con horas acordadas por fase. \\
Costo & Línea base de S/ 11,663.47 y gasto en efectivo de S/ 2,507.44. & Costos de operación del año 1: S/ 5,032.83. \\
Beneficio & Satisfacción medida en el piloto. & Beneficio de S/ 41.64 por arrendatario al año, medido desde 2027. \\
\bottomrule
\end{tabularx}
}

\begin{hojah}{Proceso de Cierre de la Fase (PMBOK)}
\begin{grande}[12.6cm]
\begin{tikzpicture}[>=Stealth]
\newcommand{\paso}[6]{
\node[draw=#4, fill=#4!10, line width=1.3pt, rounded corners=6pt, minimum width=5.6cm, minimum height=4.4cm] (#1) at (#2) {};
\node[circle, fill=#4, text=white, font=\Large\bfseries, minimum size=0.9cm] at ($(#2)+(-2.25,1.7)$) {#3};
\node[font=\large\bfseries, text=#4, anchor=west, text width=4.3cm] at ($(#2)+(-1.7,1.7)$) {#5};
\node[font=\large, text width=5.1cm, align=left] at ($(#2)+(0,-0.4)$) {#6};}
\paso{a}{0,0}{1}{AzulUNA}{Verificar\\entregables}{Cada entregable contra su criterio de aceptación. \textbf{Hoy:} 21 de 28 documentos y MVP en la nube.}
\paso{b}{6.0,0}{2}{AzulUNA}{Aceptación formal}{Votación final del sponsor y acta de cierre firmada. \textbf{Fecha:} 18/12.}
\paso{c}{12.0,0}{3}{AzulUNA}{Cerrar contratos}{Términos de los planes gratuitos revisados y cuenta de Google Play definida. \textbf{Fecha:} 11/12.}
\paso{d}{18.0,0}{4}{AzulUNA}{Liberar recursos}{Horas del semestre, celular de pruebas y accesos de cada integrante. \textbf{Fecha:} 18/12.}
\paso{e}{18.0,-5.2}{5}{Rama4}{Archivar}{Líneas base, cambios, actas e informes en el repositorio con su versión final. \textbf{Fecha:} 16/12.}
\paso{f}{12.0,-5.2}{6}{Rama4}{Lecciones\\aprendidas}{Lo que funcionó, lo que no y qué haremos distinto, por rol. \textbf{Primer bosquejo:} hoy.}
\paso{g}{6.0,-5.2}{7}{Rama4}{Comunicar\\el cierre}{Informe final y exposición ante el asesor y los clientes del piloto. \textbf{Fecha:} 14/12.}
\paso{h}{0,-5.2}{8}{Rama4}{Medir beneficios}{Satisfacción del piloto ahora y beneficio por arrendatario desde 2027.}
\foreach \x/\y in {a/b,b/c,c/d,e/f,f/g,g/h} {\draw[->, line width=2pt, AzulUNA] (\x) -- (\y);}
\draw[->, line width=2pt, AzulUNA] (d.south) -- (e.north);
\end{tikzpicture}
\end{grande}
\vspace{0.2cm}
{\large\noindent Los pasos 1 a 4 cierran lo administrativo; los pasos 5 a 8 conservan el conocimiento y el valor para las fases siguientes.}
\end{hojah}

\section{Lista de Verificación del Cierre}
\newcommand{\ok}{\textcolor{Rama4}{\large$\checkmark$}}
\newcommand{\pend}{\textcolor{GrisOscuro}{\large$\square$}}
{\small
\rowcolors{2}{white}{GrisFila}
\noindent\begin{tabularx}{\textwidth}{@{}L P{3.0cm} C{1.4cm} C{2.0cm}@{}}
\toprule
\rowcolor{AzulClaro}
\textbf{Condición} & \textbf{Responsable} & \textbf{Fecha} & \textbf{Cumplido} \\
\midrule
MVP de la Unidad I entregado y sustentado. & Eliseo Tarqui Ajahuana & 09/10 & \marca \\
Plan de Dirección aprobado y documentos de planificación en el repositorio. & Clyde N. Paricahua Pari & 26/10 & \marca \\
MVP publicado en la nube y probado con cuartos de compañeros. & Eliseo Tarqui Ajahuana & 13/11 & \marca \\
Registro e inicio de sesión en la nube funcionando, y cada usuario ve solo sus datos. & Eliseo Tarqui Ajahuana & 18/11 & \marca \\
Piloto con 10 arrendatarios y 5 propietarios, con encuesta de satisfacción. & Clyde N. Paricahua Pari & 27/11 & \marca \\
0 errores críticos abiertos y 14 casos de prueba aprobados. & Eliseo Tarqui Ajahuana & 04/12 & \marca \\
Manuales de usuario y técnico, y versión 1.0 etiquetada. & Eliseo Tarqui Ajahuana & 04/12 & \marca \\
Último corte de valor ganado y costo final. & Ruth K. Apaza Solis & 11/12 & \marca \\
Lecciones aprendidas de todos los roles completas. & Viviana M. Quisbert Quispe & 11/12 & \marca \\
Informe final expuesto. & Clyde N. Paricahua Pari & 14/12 & \marca \\
Acta de cierre firmada y repositorio archivado. & Clyde N. Paricahua Pari & 18/12 & \marca \\
\bottomrule
\end{tabularx}
}

\Needspace{18\baselineskip}
\section{Entregables y Criterios de Aceptación}
{\small
\rowcolors{2}{white}{GrisFila}
\noindent\begin{tabularx}{\textwidth}{@{}C{0.6cm} P{3.0cm} L C{1.6cm} C{2.0cm}@{}}
\toprule
\rowcolor{AzulClaro}
\textbf{N.\textsuperscript{o}} & \textbf{Entregable} & \textbf{Criterio de aceptación} & \textbf{Estado} & \textbf{Cumplido} \\
\midrule
1 & Expediente de gestión & Documentos de la ruta revisados, con cifras coherentes y en el repositorio. & \etiqueta{SevMedia}{21 de 28} & \marca \\
2 & Prototipo de la Unidad I & Ranking explicado de 5 cuartos en el orden esperado. & \etiqueta{Rama4}{Aceptado} & \marca \\
3 & MVP final & Las 9 funciones básicas aceptadas y 0 errores graves. & \etiqueta{SevMedia}{En curso} & \marca \\
4 & Piloto con clientes & 10 arrendatarios y 5 propietarios; 80\,\% o más satisfecho. & \etiqueta{GrisOscuro}{27/11} & \marca \\
5 & Cierre & Informe final, lecciones aprendidas y acta de cierre firmada. & \etiqueta{GrisOscuro}{18/12} & \marca \\
\bottomrule
\end{tabularx}
}

\vspace{0.1cm}
{\small\noindent\textbf{Pasa a la fase 2 (no es incumplimiento):} las funciones nuevas (mapa de cuartos, calificaciones, reportes, cobros y nuevas ciudades) y la publicación en Google Play. Quedan en la lista de pendientes del producto con su prioridad y se registran como exclusiones en el acta de cierre.}

\begin{hojah}{Ciclo de Vida del Software: Despliegue y Pruebas Continuas en Producción}
\begin{grande}[11.8cm]
\begin{tikzpicture}[>=Stealth]
\newcommand{\et}[5]{\node[draw=#3, fill=#3!#5, line width=1.3pt, rounded corners=6pt, minimum width=3.6cm, minimum height=2.6cm, align=center, font=\large\bfseries, text=#4] (#1) at (#2) {};}
\et{r}{0,0}{GrisOscuro}{black}{10}
\et{d}{4.2,0}{GrisOscuro}{black}{10}
\et{c}{8.4,0}{GrisOscuro}{black}{10}
\et{p}{12.6,0}{GrisOscuro}{black}{10}
\et{dep}{16.8,0}{RaciR}{RaciR}{12}
\et{pc}{21.0,0}{RaciR}{RaciR}{12}
\foreach \n/\t/\s in {r/Requisitos/9 funciones básicas, d/Diseño/{arquitectura\\y pantallas}, c/Construcción/sprints semanales, p/Pruebas/{12 automáticas\\y 14 casos}, dep/Despliegue/MVP en la nube, pc/{Pruebas\\continuas}/en producción} {
  \node[font=\large\bfseries, align=center] at ($(\n)+(0,0.5)$) {\t};
  \node[font=\normalsize, align=center] at ($(\n)+(0,-0.6)$) {\s};}
\foreach \a/\b in {r/d,d/c,c/p,p/dep,dep/pc} {\draw[->, line width=2pt, AzulUNA] (\a) -- (\b);}
\draw[->, line width=2pt, Rama4, dashed] (pc.south) -- ++(0,-0.9) -| node[pos=0.25, below, font=\large\bfseries, text=Rama4]{opinión de los usuarios a la lista de pendientes} (c.south);
\node[font=\Large\bfseries, text=RaciR] at (18.9,1.75) {Etapas de esta fase};
\draw[RaciR, line width=1.2pt, dashed, rounded corners=8pt] (14.8,-1.4) rectangle (23.0,2.15);
% detalle despliegue
\node[draw=RaciR, fill=white, line width=1.1pt, rounded corners=5pt, text width=10.6cm, align=left, font=\large, anchor=north west] at (0,-3.3) {\textbf{\textcolor{RaciR}{Despliegue}}\\[2pt]
1. Cada cambio del código lo revisa otro integrante antes de juntarlo con la versión principal.\\
2. Pasan las 12 pruebas automáticas y el programa se arma sin errores.\\
3. Se publica la nueva versión en la nube con su número (hasta la versión 1.0).\\
4. Prueba rápida: ingresar, ver el ranking, abrir un cuarto y publicar.\\
5. Si algo falla, se vuelve a la versión anterior en el mismo día.};
\node[draw=RaciR, fill=white, line width=1.1pt, rounded corners=5pt, text width=11.4cm, align=left, font=\large, anchor=north west] at (11.4,-3.3) {\textbf{\textcolor{RaciR}{Pruebas continuas en producción}}\\[2pt]
1. Cuartos reales de compañeros publicados y revisados cada semana.\\
2. Ranking desde distintos puntos de Puno: orden y desglose correctos.\\
3. Tiempo de respuesta del ranking menor a 250 ms.\\
4. Uso de los servicios gratuitos por debajo del 70\,\% de lo permitido.\\
5. Cada error encontrado se anota como tarjeta en el tablero.};
\end{tikzpicture}
\end{grande}
\end{hojah}

\section{Cierre Técnico y Arquitectura}
El cierre técnico asegura que el MVP quede operable y mantenible para las fases siguientes, y no solo entregado.

\vspace{0.1cm}
{\small
\rowcolors{2}{white}{GrisFila}
\noindent\begin{tabularx}{\textwidth}{@{}P{3.2cm} L C{1.6cm} C{2.0cm}@{}}
\toprule
\rowcolor{AzulClaro}
\textbf{Evidencia técnica} & \textbf{Qué debe existir antes de firmar el cierre} & \textbf{Estado} & \textbf{Cumplido} \\
\midrule
Documentación técnica & Cómo está hecho el sistema, sus 8 tablas de datos y la guía de instalación, todo en el repositorio. & Hecho & \marca \\
Manual de operación & Cómo publicar una versión, volver a la anterior, restablecer datos y revisar los límites de los servicios. & 04/12 & \marca \\
Transición a operaciones & Responsables del año 1 definidos y accesos a la nube transferidos a la cuenta del equipo. & 18/12 & \marca \\
Seguridad y cumplimiento & Cada usuario ve solo sus datos, contraseñas protegidas, copia de seguridad semanal y consentimiento según la Ley de Protección de Datos Personales. & 04/12 & \marca \\
Pendientes técnicos ordenados & Lista de mejoras no hechas, con su prioridad, para la fase 2. & 11/12 & \marca \\
Métricas de sostenibilidad & Costo mensual, tiempo de respuesta y disponibilidad del MVP medidos después del cierre. & Desde 18/12 & \marca \\
\bottomrule
\end{tabularx}
}

\Needspace{14\baselineskip}
\subsection{Pendientes técnicos conocidos}
{\small
\rowcolors{2}{white}{GrisFila}
\noindent\begin{tabularx}{\textwidth}{@{}L C{1.6cm} P{3.2cm}@{}}
\toprule
\rowcolor{AzulClaro}
\textbf{Pendiente} & \textbf{Prioridad} & \textbf{Destino} \\
\midrule
El puntaje llega al tope de 100 en el mejor cuarto de todos los perfiles; revisar el margen de las bonificaciones. & \etiqueta{SevAlta}{Alta} & Antes de la v1.0 \\
La app Android aún usa la base local; falta conectarla a la misma nube que la web. & \etiqueta{SevAlta}{Alta} & Fase 2 \\
Mapa con todos los cuartos, confirmación de alquiler y calificaciones. & \etiqueta{SevMedia}{Media} & Fase 2 \\
Reportes de publicaciones falsas y panel de administración. & \etiqueta{SevMedia}{Media} & Fase 2 \\
Pruebas automáticas del servidor y de las pantallas, además de las del puntaje. & \etiqueta{SevMedia}{Media} & Fase 2 \\
\bottomrule
\end{tabularx}
}

\Needspace{20\baselineskip}
\section{Revisión de Riesgos Cada Dos Semanas}
Desde la semana 12 el registro de riesgos se revisa a fondo cada dos semanas, además del control semanal de los disparadores. En cada revisión se recalifican todos los riesgos, se cierran los que pasaron su ventana y se agregan los del cierre.

\vspace{0.1cm}
\begin{center}
\begin{tikzpicture}[font=\small]
\draw[line width=1.5pt, AzulUNA] (0,0) -- (15.6,0);
\foreach \x/\d/\t in {0.6/13-11/{Revisión 1\\riesgos del cierre}, 4.6/27-11/{Revisión 2\\tras el piloto}, 8.6/11-12/{Revisión 3\\antes del informe}, 12.6/18-12/{Revisión final\\riesgos al año 1}} {
  \node[circle, fill=RaciR, minimum size=0.35cm] at (\x,0) {};
  \node[above=4pt, font=\small\bfseries, text=RaciR] at (\x,0.15) {\d};
  \node[below=6pt, align=center] at (\x,-0.1) {\t};}
\end{tikzpicture}
\end{center}

{\small
\rowcolors{2}{white}{GrisFila}
\noindent\begin{tabularx}{\textwidth}{@{}L C{0.6cm} C{0.6cm} C{0.9cm} P{4.2cm}@{}}
\toprule
\rowcolor{AzulClaro}
\textbf{Riesgo del cierre} & \textbf{P} & \textbf{I} & \textbf{Sev.} & \textbf{Respuesta} \\
\midrule
El registro en la nube no queda listo para el piloto. & 3 & 3 & \etiqueta{SevAlta}{\,9\,} & Ruth apoya a Eliseo y el piloto pasa a la semana 14. \\
Los clientes del piloto no aceptan el MVP como útil. & 2 & 3 & \etiqueta{SevAlta}{\,6\,} & Guion de prueba y encuesta con criterios claros. \\
El MVP en la nube se cae o supera los límites gratuitos. & 2 & 2 & \etiqueta{SevMedia}{\,4\,} & Monitoreo semanal y vuelta a la versión anterior. \\
Las lecciones aprendidas quedan incompletas por los exámenes. & 2 & 2 & \etiqueta{SevMedia}{\,4\,} & Plantilla por rol entregada el 13/11. \\
No se define quién opera el producto en el año 1. & 2 & 3 & \etiqueta{SevAlta}{\,6\,} & Acuerdo de continuidad en el acta de cierre. \\
\bottomrule
\end{tabularx}
}

\Needspace{24\baselineskip}
\section{Cálculos de Cierre (Preliminares)}
Con el corte del 06/11 y el pronóstico al terminar. Se recalculan con el último corte del 11/12.

\vspace{0.1cm}
{\small
\rowcolors{2}{white}{GrisFila}
\noindent\begin{tabularx}{\textwidth}{@{}P{3.4cm} P{3.0cm} Q{2.4cm} L@{}}
\toprule
\rowcolor{AzulClaro}
\textbf{Indicador} & \textbf{Fórmula} & \textbf{Al 06/11} & \textbf{Lectura para el cierre} \\
\midrule
Variación del costo & CV = EV $-$ AC & $-$342.68 & Se gastó un poco más de lo ganado. \\
Índice del costo & CPI = EV / AC & 0.95 & En el límite verde; meta final 0.95 o más. \\
Variación del cronograma & SV = EV $-$ PV & $-$958.73 & Cerca de una semana de atraso en la nube. \\
Índice del cronograma & SPI = EV / PV & 0.87 & Debe llegar a 1.00 si se completa el alcance de la fase 1. \\
Costo final pronosticado & EAC = BAC / CPI & 11,700.40 & S/ 36.93 sobre la línea base; se cubre con la contingencia. \\
Varianza de presupuesto & BAC $-$ EAC & $-$592.33 & Pronóstico ligeramente sobre el presupuesto del trabajo. \\
Realización de beneficios & real / plan $\times$ 100 & -- & Se mide desde 2027 (plan del año 1: S/ 14,574.00). \\
\bottomrule
\end{tabularx}
}

\Needspace{16\baselineskip}
\subsection{Criterios para declarar el resultado}
{\small
\rowcolors{2}{white}{GrisFila}
\noindent\begin{tabularx}{\textwidth}{@{}P{2.0cm} L C{2.2cm} C{2.0cm}@{}}
\toprule
\rowcolor{AzulClaro}
\textbf{Criterio} & \textbf{Se considera cumplido si} & \textbf{Tendencia hoy} & \textbf{Cumplido} \\
\midrule
Alcance & Las 9 funciones básicas aceptadas y lo que pasa a la fase 2 queda escrito en el acta. & \etiqueta{SevMedia}{6 de 9} & \marca \\
Tiempo & El acta de cierre se firma dentro de la semana 17. & \etiqueta{Rama4}{En plazo} & \marca \\
Costo & CPI final de 0.95 o más y gasto en efectivo de S/ 2,507.44 o menos. & \etiqueta{Rama4}{CPI 0.95} & \marca \\
Valor & 80\,\% o más de satisfacción en el piloto y MVP operando en la nube. & \etiqueta{SevMedia}{Por medir} & \marca \\
\bottomrule
\end{tabularx}
}

\vspace{0.1cm}
{\small\noindent Con estos criterios, el resultado de la fase 1 será \textbf{éxito} si se cumplen los cuatro, y \textbf{éxito parcial} si el alcance diferido está formalizado y el piloto confirma el valor del MVP.}

\Needspace{20\baselineskip}
\section{Lecciones Aprendidas Preliminares}
Cada lección se escribe con cuatro partes: contexto, qué ocurrió, impacto y qué haremos distinto. Se reutilizan en las fases siguientes del producto.

\vspace{0.1cm}
{\footnotesize
\setlength{\tabcolsep}{3pt}
\rowcolors{2}{white}{GrisFila}
\noindent\begin{tabularx}{\textwidth}{@{}C{0.7cm} P{2.6cm} L P{2.5cm} P{3.6cm}@{}}
\toprule
\rowcolor{AzulClaro}
\textbf{N.\textsuperscript{o}} & \textbf{Contexto} & \textbf{Qué ocurrió} & \textbf{Impacto} & \textbf{Qué haremos distinto} \\
\midrule
\multicolumn{5}{@{}l}{\cellcolor{Rama4!15}\textbf{Lo que funcionó}} \\
L1 & Construcción (sem. 4 a 8) & El motor se separó de las pantallas y se probó con pruebas automáticas. & MVP de la Unidad I 3 días antes. & Separar la lógica desde el primer sprint. \\
L2 & Seguimiento (todo el proyecto) & Tablero compartido con evidencia obligatoria para pasar a Hecho. & Control claro de quién hizo qué. & Mantener la regla en la operación. \\
L3 & Revisión del sprint (viernes) & Demostración semanal del incremento al equipo como sponsor. & Defectos detectados en días, no en semanas. & Invitar a clientes reales a una demostración por mes. \\
\multicolumn{5}{@{}l}{\cellcolor{SevAlta!12}\textbf{Lo que debemos mejorar}} \\
L4 & Planificación (sem. 4 a 10) & Los documentos de planificación se firmaron con dos semanas de retraso. & En la semana 6 se había avanzado solo el 56\,\% de lo planeado. & Planificación mínima firmada antes de construir. \\
L5 & Nube (sem. 10 y 11) & No se probó antes cómo dar permisos por usuario en la nube. & El registro se atrasó 6 días. & Probar cada herramienta nueva en la primera semana de uso. \\
L6 & Campo (sem. 6 a 12) & La meta de 15 cuartos se fijó sin medir el ritmo de las visitas. & Meta reducida a 12 con un cambio. & Estimar metas de campo con una semana de prueba. \\
L7 & Equipo (todo el proyecto) & La programación recayó casi toda en una sola persona. & Si él se detenía, todo se detenía. & Rotar tareas sencillas de código entre todos. \\
L8 & Riesgos (todo el proyecto) & Cada martes se revisaban solo las señales de alerta de los riesgos. & El problema de la nube se vio tarde. & Revisión completa del registro cada dos semanas. \\
\bottomrule
\end{tabularx}
}

\newcommand{\espacio}{\rule{0pt}{1.25cm}}
\newcommand{\perspectiva}[4]{%
\Needspace{18\baselineskip}
\subsection{#1}
{\small\noindent\textbf{Quién la completa:} #2 \quad \textbar\quad \textbf{Preguntas guía:} \textit{#3}}

\vspace{0.1cm}
{\small
\noindent\begin{tabularx}{\textwidth}{@{}P{3.6cm} |L@{}}
\toprule
\rowcolor{AzulClaro}
\textbf{Aspecto} & \textbf{Respuesta} \\
\midrule
Lo que salió bien & \espacio \\ \hline
Lo que salió mal & \espacio \\ \hline
Qué haría distinto & \espacio \\ \hline
#4 & \espacio \\
\bottomrule
\end{tabularx}
}
\vspace{0.15cm}}

\newcommand{\perspectivallena}[6]{%
\Needspace{16\baselineskip}
\subsection{#1}
{\small\noindent\textbf{Responde:} #2}

\vspace{0.1cm}
{\small
\rowcolors{2}{white}{GrisFila}
\noindent\begin{tabularx}{\textwidth}{@{}P{3.0cm} L@{}}
\toprule
\rowcolor{AzulClaro}
\textbf{Aspecto} & \textbf{Respuesta} \\
\midrule
Lo que salió bien & #3 \\
Lo que salió mal & #4 \\
Qué haría distinto & #5 \\
Mirando adelante & #6 \\
\bottomrule
\end{tabularx}
}
\vspace{0.15cm}}
\section{Perspectivas del Cierre por Rol}
Cada integrante respondió con hechos concretos. Los cuadros de los clientes, del sponsor y del asesor se completan después del piloto. Estas respuestas alimentan el informe final y las lecciones aprendidas definitivas.

\perspectivallena{Líder del proyecto}{Clyde Neil Paricahua Pari, Project Manager}{Lo que más me sirvió fue trabajar por semanas, con una revisión cada viernes y un tablero en GitHub donde ninguna tarea se daba por terminada sin su evidencia. Así siempre supe quién hizo qué y qué faltaba. También creo que acerté al dejar Horus Sentinel el 22/09: fue la decisión más difícil, porque ya llevábamos 4 semanas en él, pero los números mostraban que iba atrasado y que daba mucho menos beneficio que CasaScore. Dejarlo nos permitió enfocarnos en un solo proyecto.}{Nos demoramos en planificar: los documentos se firmaron con dos semanas de retraso y en la semana 6 solo habíamos avanzado el 56\,\% de lo previsto. Al inicio no sabíamos cuánto debía durar realmente el proyecto; solo que debía llegar a ser rentable. La programación recayó en Eliseo, y la meta de 15 cuartos la fijé sin medir cuántos podíamos visitar, así que la bajamos a 12. Además, no todos participaban igual en las reuniones y muchas veces tuve que insistir o asumir tareas.}{Planificaría antes de construir y fijaría desde el primer día cuánto dura cada fase y qué incluye, con un tiempo exacto para cada tarea. Revisaría todos los riesgos cada dos semanas desde el inicio. Buscaría que nos organicemos mejor como grupo, que lleguemos sabiendo más de las herramientas y delegaría más desde el comienzo.}{En los primeros 3 años la app no recupera lo invertido con su propio dinero, así que el reto es mantenerla sin gastar más de lo que genera. Según el análisis a 10 años, quedándonos en la región Puno la inversión se recupera en unos 3.6 años. Hay que definir quién la mantiene después del semestre y postular a concursos y fondos. En lo personal, aprendí a decidir con datos y no por intuición, a delegar y a usar Git y GitHub.}
\perspectivallena{Programador}{Eliseo Tarqui Ajahuana, Especialista Técnico}{Separar desde la semana 7 la parte que calcula el puntaje de las pantallas nos salvó: la probamos con pruebas automáticas y el primer prototipo salió 3 días antes. La forma de publicar también funcionó: cada cambio revisado por otro compañero, las 12 pruebas pasando y la app quedó en la nube el 13/11, probada con cuartos reales de compañeros.}{No probamos a tiempo cómo dar permisos por usuario en la nube, y por eso el registro se atrasó 6 días. Casi toda la programación cayó en mí, así que cualquier problema me frenaba solo a mí. Además, la app del celular sigue guardando los datos en el propio teléfono y aún no se conecta a la nube; eso queda como tarea importante para la fase 2.}{Haría una prueba corta de los permisos en la nube en la primera semana, no después. Repartiría tareas sencillas de programación entre todos para no ser el único que avanza. Y dejaría listas desde antes las pruebas automáticas del servidor y de las pantallas, no para el final.}{La app quedó lista para seguir usándose y mejorándose: está documentado cómo está hecha, sus 8 tablas de datos y cómo instalarla. Para la fase 2 priorizaría conectar la app del celular a la nube, revisar por qué el mejor cuarto siempre llega a 100 puntos y agregar más pruebas automáticas, para que el año 1 empiece sin problemas graves.}
\perspectivallena{Analista de planificación}{Ruth Karina Apaza Solis}{El seguimiento con valor ganado mostró el atraso con números y a tiempo: pasamos de avanzar el 56\,\% de lo planeado en la semana 6 al 87\,\% al 06/11. El costo se mantuvo controlado: por cada sol gastado obtuvimos 95 céntimos de trabajo, y el costo final previsto (S/ 11,700.40) solo supera el plan en S/ 36.93, que cubre el dinero reservado para riesgos. Las tarifas por hora del equipo se respetaron.}{La parte de la nube se estimó sin probarla antes, y cuando falló tuve que apoyar a Eliseo para sacarla adelante. Los primeros cortes se hicieron sin un plan aprobado, así que esos números sirvieron más como alerta que como medida exacta.}{Haría el corte de avance cada viernes con los datos del tablero, y no después de varias semanas. Antes de estimar una herramienta nueva, haría una prueba corta. Y probaría una semana antes de fijar metas de campo, como la de los cuartos.}{Propongo comparar cada trimestre desde 2027 el beneficio real con el previsto para el año 1 (S/ 14,574.00), y revisar cada mes que el uso de los servicios gratuitos en la nube no pase del 70\,\% de lo permitido. Las funciones de la fase 2 deben estimarse con la velocidad real que tuvimos en esta fase.}
\perspectivallena{Gestora de comunicación, riesgos y documentación}{Viviana Mary Quisbert Quispe}{Las 5 actas de las semanas 4 a 8 registraron el avance, los acuerdos y los compromisos con su responsable, y sus cifras coinciden con los informes de control. La señal de alerta del riesgo de la nube detectó a tiempo que el inicio de sesión no funcionaba. La convocatoria al piloto logró 14 inscritos para una meta de 10, y así bajó el riesgo de no tener participantes.}{Mis documentos (políticas, riesgos, comunicaciones y actas) llegaron al repositorio recién el 06/11, aunque la regla del equipo pedía subir cada acta en 24 horas. Como solo revisaba las señales de alerta y no todos los riesgos, el problema de la nube se vio tarde: el registro se atrasó 6 días y el piloto pasó a la semana 14.}{Escribiría el acta durante la reunión y la subiría ese mismo día. Pondría cada riesgo como tarjeta en el tablero desde el primer día, para que todo el equipo lo vea. Y para lo técnico nuevo pondría alertas más tempranas, por ejemplo probar la nube en la semana 8 y no en la 10.}{Para el año 1 los riesgos principales serán que la app se caiga o supere lo que permiten los servicios gratuitos, que nadie quede a cargo de mantenerla, las publicaciones falsas y la protección de datos personales. Propongo revisar los riesgos cada mes, hacer una copia de seguridad semanal y usar mis 2 horas semanales de moderación para revisar los reportes.}
\perspectiva{Cliente arrendatario}{Participantes del piloto (encuesta y entrevista)}{facilidad para encontrar cuarto, si el puntaje se entiende y si confían en el ranking}{Si recomendaría la app y por qué}
\perspectiva{Cliente propietario}{Propietarios del piloto (entrevista)}{tiempo para publicar, fotos obligatorias, interés recibido de inquilinos}{Si pagaría S/ 5 al mes por cuarto publicado}
\perspectivallena{Sponsor}{Team Sobrevivientes en votación (los cuatro integrantes)}{Todo el trabajo quedó documentado: cada documento está en el repositorio y cada tarea tiene su evidencia en el tablero, así que cualquiera puede revisar qué se hizo, quién lo hizo y cuándo. También funcionó el diálogo: las decisiones importantes, como dejar Horus Sentinel o mover el piloto, se conversaron con datos en las reuniones de los martes y viernes.}{Hubo atrasos: la planificación se firmó dos semanas tarde y la parte de la nube se demoró, por lo que el piloto pasó a la semana 14. Además, la carga no se repartió igual: la programación recayó casi toda en una persona y algunas tareas se concentraron en pocos integrantes.}{No aprobaríamos empezar a construir sin una planificación mínima firmada. Revisaríamos los indicadores cada semana y no solo al final de cada fase. Y fijaríamos desde el inicio cuánto debe durar cada fase y en qué momento la app debe empezar a ser rentable.}{La decisión sobre continuar con la fase 2 se votará después del piloto (23 al 27/11), con la satisfacción de los clientes y el estado del MVP a la vista.}
\perspectiva{Asesor}{Ing. Renzo Apaza Cutipa (opcional)}{calidad del expediente, aplicación de los marcos de referencia, observaciones recurrentes}{Recomendación para el equipo}

\Needspace{20\baselineskip}
\section{Bosquejo del Informe Final}
{\small
\rowcolors{2}{white}{GrisFila}
\noindent\begin{tabularx}{\textwidth}{@{}C{0.7cm} P{4.2cm} L P{2.5cm}@{}}
\toprule
\rowcolor{AzulClaro}
\textbf{N.\textsuperscript{o}} & \textbf{Sección} & \textbf{Contenido} & \textbf{Responsable} \\
\midrule
1 & Resumen ejecutivo & Qué se logró, resultado del piloto y decisión de éxito. & Clyde \\
2 & Objetivos y resultados & O1 a O5 con su meta y su valor final. & Clyde \\
3 & Producto entregado & MVP en la nube, requisitos aceptados y exclusiones a la fase 2. & Eliseo \\
4 & Desempeño & Avance contra el plan, costo final, curva de avance y fechas cumplidas. & Ruth \\
5 & Cambios y riesgos & Solicitudes de cambio y riesgos materializados con su respuesta. & Viviana \\
6 & Perspectivas por rol & Resumen de los cuadros de cada participante. & Todos \\
7 & Lecciones aprendidas & Lecciones definitivas con su acción para el futuro. & Viviana \\
8 & Continuidad & Hoja de ruta 2027 a 2029 y responsables de la operación. & Clyde \\
\bottomrule
\end{tabularx}
}

\Needspace{18\baselineskip}
\section{Bosquejo del Acta de Cierre}
{\small
\noindent\begin{tabularx}{\textwidth}{@{}P{3.8cm} L@{}}
\toprule
\textbf{Proyecto y fase:} & CasaScore (AlquilaPuno), fase 1: proyecto del semestre 2026-II. \\
\textbf{Entregables aceptados:} & Expediente de gestión, prototipo de la Unidad I, MVP final en la nube, piloto con clientes e informe final. \\
\textbf{Exclusiones:} & Mapa de cuartos, calificaciones, reportes, cobros dentro de la app, nuevas ciudades y Google Play, que pasan a la fase 2. \\
\textbf{Resultados:} & Avance final, costo final, satisfacción del piloto y decisión de éxito (se completan el 18/12). \\
\textbf{Continuidad:} & Responsables de la operación del año 1 y fecha de inicio de la fase 2. \\
\textbf{Firmas:} & Los cuatro integrantes como sponsor y equipo del proyecto. \\
\bottomrule
\end{tabularx}
}

\begin{hojah}{Perspectiva del Producto: Hoja de Ruta a Tres Años}
\begin{grande}[12.4cm]
\begin{tikzpicture}[>=Stealth]
\draw[line width=3pt, AzulUNA, ->] (0,0) -- (24.2,0);
\foreach \x/\t/\c/\s/\d in {0.3/2026-II: fase 1/RaciR/{Proyecto del semestre}/{MVP en la nube, piloto con compañeros y clientes, expediente y cierre}, 6.3/2027: año 1/Rama4/{Operación en Puno}/{500 usuarios al mes, app en Google Play, suscripción de S/ 5 por cuarto y anuncios}, 12.3/2028: año 2/Rama2/{Crecimiento}/{1,300 usuarios al mes, llegada a Juliaca, calificaciones y reportes de publicaciones}, 18.3/2029: año 3/Rama5/{Expansión}/{2,100 usuarios al mes en la región Puno; desde 2030, otras ciudades si se cumplen las metas}} {
  \node[circle, fill=\c, minimum size=0.6cm] at (\x+2.6,0) {};
  \node[fill=\c, text=white, font=\large\bfseries, minimum width=5.4cm, minimum height=0.9cm, anchor=south west] at (\x,0.6) {\t};
  \node[draw=\c, fill=\c!8, line width=1.2pt, text width=5.0cm, align=left, font=\normalsize, anchor=north west, minimum height=4.0cm] at (\x,-0.6) {\textbf{\s}\\[3pt]\d};}
\node[font=\large\bfseries, text=RaciR, align=center] at (2.9,3.0) {Se cierra en diciembre};
\draw[RaciR, line width=1.4pt, dashed] (5.95,-5.0) -- (5.95,3.3);
\node[font=\normalsize, text width=23cm, align=left, anchor=north west] at (0,-5.3) {Beneficio esperado para los arrendatarios: S/ 14,574 (año 1), S/ 37,892 (año 2) y S/ 61,211 (año 3). Con su propio dinero, la app recupera la inversión en unos 3.6 años si se queda en la región Puno.};
\end{tikzpicture}
\end{grande}
\end{hojah}

\Needspace{24\baselineskip}
\section{Firmas}
\noindent En señal de conformidad con el presente bosquejo, los integrantes del equipo firman a continuación:

\vspace{1.5cm}
\begin{center}
\begin{tabular}{c c c}
    \parbox{6.5cm}{\centering \rule{6.0cm}{0.6pt}\\[5pt] \textbf{Clyde Neil Paricahua Pari}\\[2pt] {\footnotesize Project Manager}}
    & \hspace{0.8cm} &
    \parbox{6.5cm}{\centering \rule{6.0cm}{0.6pt}\\[5pt] \textbf{Ruth Karina Apaza Solis}\\[2pt] {\footnotesize Analista de Planificación}}
    \\[2.0cm]
    \parbox{6.5cm}{\centering \rule{6.0cm}{0.6pt}\\[5pt] \textbf{Eliseo Tarqui Ajahuana}\\[2pt] {\footnotesize Especialista Técnico / Arquitectura}}
    & \hspace{0.8cm} &
    \parbox{6.5cm}{\centering \rule{6.0cm}{0.6pt}\\[5pt] \textbf{Viviana Mary Quisbert Quispe}\\[2pt] {\footnotesize Gestora de Comunicación, Riesgos y Documentación}}
\end{tabular}
\end{center}


\end{document}
